\chapter{Validation And Oracles}

\section{Plural Oracles}

There is no single oracle for APFS indexing. The correct oracle depends on the
metric and product mode:

\begin{itemize}
  \item namespace: POSIX/API traversal of the same selected volume or stable view
  \item logical size: public file metadata APIs and raw inode/dstream comparison
  \item allocated size: metric-specific public tool and raw extent evidence
  \item incremental correctness: fresh full raw parse of the same selected state
  \item boot-root semantics: user-visible macOS namespace only when that mode is
  explicitly requested
\end{itemize}

The canonical oracle policy is in
\path{docs/research/RL-10-validation-corpus-and-oracle.md}.

\section{Successful Proof Pattern}

The successful v1 proof pattern is:

\begin{enumerate}
  \item mount a lab image
  \item build a deterministic corpus
  \item capture mounted oracle
  \item detach the image to stabilize raw state
  \item reattach without mounting
  \item raw-walk the APFS container
  \item normalize and compare path, type, file identity, logical size, and
  symlink target
\end{enumerate}

This pattern was established in \path{docs/research/experiments/EX-03-pinned-state-raw-vs-oracle/README.md}
and expanded in \path{docs/research/experiments/EX-04-expanded-pinned-corpus/README.md}.

\section{Negative Live-State Pattern}

\path{docs/research/experiments/EX-05-live-pinned-churn/README.md} is the
counterexample pattern. A mounted image can be raw-readable while latest-state
raw output still lacks a named stable oracle. That result prevents the manual
from treating live latest raw traversal as correctness evidence.
